\chapter{Introduction}

\section{Background and Motivation}
Virtual voice assistants have become ubiquitous since the introduction of Apple’s Siri in 2011, offering convenient hands-free interaction through natural speech\,[1]. As these conversational agents handle an expanding range of sensitive tasks (from smart-home control to banking), ensuring their safety and security has grown paramount. Unfortunately, voice-based systems are inherently vulnerable to spoofing and impersonation attacks—malicious actors can exploit the interface by injecting falsified audio commands or masquerading as a trusted speaker. Recent advances in AI-driven voice cloning (deepfakes) have enabled attackers to generate synthetic speech that is nearly indistinguishable from a target speaker, leading to real-world incidents. For example, in 2019 fraudsters used AI to mimic a CEO’s voice and trick an employee into a \$243{,}000 bank transfer\,[2]. Likewise, researchers have demonstrated that even state-of-the-art voice authentication systems (e.g., WeChat Voiceprint, Microsoft Azure, Amazon Alexa’s voice ID) can be bypassed by cloned voices, highlighting the severity of this threat\,[3].

\section{Problem Statement and Research Gap}
Existing countermeasures typically fall into two categories: (i) user authentication and liveness checks, and (ii) deepfake detection algorithms. The former (e.g., voice biometrics, challenge–response) can be intrusive and are known to be defeasible via replayed or synthesized responses. The latter (artifact-based detection) has seen rapid progress\,[4][5] but remains brittle under distribution shift: detectors often struggle with unseen attacks, codecs, recording chains, or acoustic environments, causing error rates to climb markedly outside controlled conditions\,[6][7][8]. In interactive assistants, this mismatch undermines safety at the exact moment of deployment. Hence, there is a clear need for a proactive, layered defense that (a) offers a lightweight provenance check for assistant-originated audio and (b) preserves strong detection capability against forged user-originated audio, even in adverse channels.

\section{Objectives and Research Questions}
To address this gap, we pursue a combined approach that joins audio watermarking for \emph{assistant speech provenance} with robust deepfake/replay detection for \emph{user inputs}. The study is guided by the following research questions:

\begin{itemize}[leftmargin=1.25em]
  \item \textbf{RQ1 (Imperceptibility):} Can a TTS watermark be embedded in the assistant’s responses without degrading perceived naturalness (MOS and A/B tests)?
  \item \textbf{RQ2 (Robustness):} Does the watermark survive re-recording, common channel effects, and lossy codecs while remaining machine-detectable\,[9][10]?
  \item \textbf{RQ3 (Detection Generalization):} Under VoIP compression, noise, and room acoustics, how well does the deepfake/replay detector maintain accuracy on \emph{unseen} attacks (DET curves, EER)\,[6][13]?
  \item \textbf{RQ4 (System Integration):} What is the end-to-end effectiveness and latency of the integrated “watermark + detector” middleware in a turn-taking voice assistant pipeline?
\end{itemize}

\section{Proposed Approach (Overview)}
We add a proactive safety layer to voice-based conversational systems by combining two complementary techniques:
\begin{enumerate}[leftmargin=1.25em]
  \item \textbf{Imperceptible TTS Watermarking} of agent speech for provenance: an imperceptible, hard-to-remove watermark is embedded in every assistant utterance. Any audio claiming to originate from the assistant can be verified via a fast watermark decoder; fakes or tampered content (lacking the signature) are rejected.
  \item \textbf{Deepfake/Replay Detection} of user speech for integrity: incoming audio is screened by a channel-robust deepfake/replay detector tuned via DET curves and EER to balance security and usability\,[13].
\end{enumerate}
By validating both agent-originated and user-originated audio in real time, the middleware halts malicious injections early, exploits complementary strengths (provenance vs.\ content analysis), and keeps runtime overhead low.

\section{Scope, Assumptions, and Threat Model}
\textbf{Adversary capabilities.} We assume attackers can: (i) synthesize convincing target voices using modern TTS/VC; (ii) replay high-quality recordings; (iii) apply common signal transforms (compression, resampling, additive noise) and attempt naive watermark removal. \textbf{Out-of-scope.} We do not consider physical compromise of the user device, sensor hardware modification, or attackers with perfect knowledge of and white-box access to the exact watermarking key and embedding internals.

\section{Evaluation Criteria and Datasets (Overview)}
We evaluate the solution with: (i) objective watermark metrics (bit/payload error rates, detection rate after distortions)\,[11][12]; (ii) deepfake/replay detection metrics (DET curves, EER, FPR at fixed TPR)\,[13]; (iii) system-level latency; and (iv) subjective quality via MOS and A/B listening tests to verify imperceptibility\,[14]. Test audio covers typical degradations (VoIP vocoders, reverberation, speaker–microphone re-recording, and noise).

\section{Contributions}
\begin{itemize}[leftmargin=1.25em]
  \item \textbf{Imperceptible TTS Watermarking.} A watermarking scheme for assistant speech that is inaudible yet resilient to re-recording and typical channel distortions (noise, filtering, compression), enabling fast authenticity checks\,[11][12].
  \item \textbf{Channel-Robust Deepfake/Replay Detection.} A state-of-the-art detector trained and calibrated with diverse degraded audio (e.g., VoIP vocoders, reverberant rooms) to maintain low EER in realistic conditions, with operating points chosen via DET analysis\,[13].
  \item \textbf{Integrated Safety Middleware.} An open-source prototype that couples the encoder/decoder and detector within the assistant pipeline to verify every dialog turn before execution, blocking spoofed inputs or forged outputs.
  \item \textbf{Comprehensive Evaluation.} A targeted benchmark spanning degradation scenarios, reporting: watermark survivability, detection accuracy under channel shift, runtime latency, and perceptual impact (MOS/A–B)\,[14].
\end{itemize}

